\documentclass[11pt]{article}
\usepackage[margin=1in]{geometry}
\usepackage{amsmath, amssymb, booktabs, hyperref, xcolor, listings}
\usepackage{enumitem}
\lstset{basicstyle=\ttfamily\small, breaklines=true, frame=single}
\title{Independent Replication Report:\\
\emph{Characterizing and Mitigating Coherent Errors in a Trapped Ion Quantum Processor Using Hidden Inverses} (Majumder et al., arXiv:2205.14225)}
\author{QC-200 replication wave, run 2026-07-05}
\date{}
\begin{document}
\maketitle

\section{Paper summary}
Majumder, Yale, Morris, Lobser, Burch, Chow, Revelle, Clark, and Pooser (Quantum, 2023; arXiv:2205.14225v2) study coherent-error characterization and mitigation on the QSCOUT trapped-ion testbed at Sandia National Labs. Their central contribution is the experimental demonstration of the \textbf{Hidden Inverses} (HI) error-mitigation protocol on single- and two-qubit gates and its use inside a small VQE calculation for H$_2$.
The HI idea is that self-adjoint unitaries (like $H$ or CNOT) can be compiled either as $G=ABC$ or as $G^{\dagger}=C^{\dagger}B^{\dagger}A^{\dagger}$; in the presence of coherent errors these two decompositions leave different residual unitaries, so a compiler can strategically choose one so that neighboring-gate errors cancel.
The paper's four experimental threads are: (a) fitting a physics-focused single-qubit noise model with over-rotation $\varepsilon$, phase $\phi$, and detuning $\delta$ (Eq.~1) using the amplification circuit $[H\,X(\Theta)\,Z(\Phi)\,H^{\dagger}]^{100}$ (Fig.~5); (b) applying HI to the Hadamard and CNOT decompositions (Fig.~1--2); (c) VQE for H$_2$ with HI vs Randomized Compiling (RC) vs density-matrix purification; (d) MS-gate fidelity budget under various noise-injection regimes.

\medskip
\textbf{Provenance check.} Title and authors were verified from the fetched PDF metadata: \texttt{pdfinfo paper.pdf} returned exactly the title used above and the nine-author list. The task brief mentioned BB1/SK1/CORPSE composite pulses as the ``expected'' mitigation technique; the paper actually uses \emph{Hidden Inverses}, which is a structurally different mechanism (structural symmetry, not envelope shaping). Our reproduction targets the paper's actual technique, not the brief's rumored technique.

\section{Claims table}
\begin{tabular}{p{0.6cm} p{7.5cm} p{2.4cm} p{2.4cm}}
\toprule
\# & Claim & Testable? & Tested? \\ \midrule
C1 & HI decomposition: $H=Y(-\pi/2)X(\pi)$ vs $H^{\dagger}=X(-\pi)Y(\pi/2)$; both implement the same operation ideally, differ under coherent noise (Fig.~1, Sec.~2). & Analytical + sim & Yes \\
C2 & Noise model Eq.~(1): $X_{\text{noisy}}=\exp[-i\Omega(1+\varepsilon)t/2\,(\cos\phi\, X+\sin\phi\, Y) - i\delta t/2\,Z]$. & Analytical & Yes \\
C3 & On $[H\,X(\Theta)\,Z(\Phi)\,H^{\dagger}]^{100}$ across a $(\Theta,\Phi)$ grid, non-linear least squares recovers $(\varepsilon,\phi,\delta)$ with variance $\sim 10^{-4}$ (Sec.~4.3). & Numerical & Yes \\
C4 & With coherent errors reversed between adjacent self-adjoint gates, HI cancels the leading order coherent error. & Numerical & Yes \\
C5 & MS-gate best-fit noiseless-baseline model: over-rot $=0.09$~rad, $p_{2q}=0.02$, $\bar n=0.05$ $\Rightarrow$ $F_{\text{MS}}\approx 97.5\%$. & Numerical & Yes \\
C6 & Under-rotation $0.45$~rad case: $F_{\text{MS}}\approx 91\%$. & Numerical & Yes \\
C7 & Reduced-cooling case: over-rot $0.12$~rad, $p_{2q}=0.06$, $\bar n=0.5$ $\Rightarrow$ $F_{\text{MS}}\approx 89\%$. & Numerical + hardware & Partly (no Debye-Waller stochastic model added) \\
C8 & VQE H$_2$ raw energy expectation values improve from $0.923$ to $0.950$ with HI, and to $0.997$ with HI+purification (Table in Sec.~5.3, ``None'' row). & Numerical + hardware & Not reproduced (VQE end-to-end out of scope for a single-shot reproduction) \\ \bottomrule
\end{tabular}

\section{Method}
All work was done on macOS/Darwin 25.3.0 x86\_64 with system Python 3 (\texttt{numpy 2.4.3}, \texttt{scipy 1.18.0}); no Qiskit/Cirq/PennyLane installation was needed. Every simulation is a direct numpy statevector or density-matrix computation.

\begin{enumerate}[leftmargin=1.4em]
\item \textbf{Fetch + verify.}
\begin{lstlisting}
curl -sL https://arxiv.org/pdf/2205.14225 -o paper.pdf
pdfinfo paper.pdf   # title/authors/11pp
pdftotext paper.pdf work/paper.txt
\end{lstlisting}
\item \textbf{Noise-model implementation} (\texttt{report/evidence/sim\_hidden\_inverses.py}, \texttt{X\_noisy, Y\_noisy}) — direct matrix exponential of the Hamiltonian in Eq.~(1); scaling-and-squaring Taylor expansion (order 30) avoids the need for a full linear-algebra library beyond numpy.
\item \textbf{H decomposition} — \texttt{H\_standard = Y(-\(\pi\)/2)@X(\(\pi\))}, \texttt{H\_hidden = X(-\(\pi\))@Y(\(\pi\)/2)} exactly as in Fig.~1.
\item \textbf{R2 (parameter fit).} 11\(\times\)11 \((\Theta,\Phi)\) grid on \([-0.08,0.08]\)~rad (matching the visible axis of Fig.~5), 100 repetitions per block, injected true parameters \((\varepsilon,\phi,\delta)=(0.03,-0.02,0.015)\) plus Gaussian shot-noise \(\sigma=0.005\), fit with \texttt{scipy.optimize.curve\_fit} (the paper explicitly names ``SciPy's non-linear least-square fitting algorithm'').
\item \textbf{R3 (HI cancellation).} \texttt{report/evidence/sim\_hi\_cancellation.py} scans over-rotation \(\varepsilon\in[-0.1,0.1]\) and computes \(1-F_{\text{avg}}\) for \(H\cdot H\) (bare) vs \(H\cdot H^{\dagger}\) (HI). Slope of \(\log(1-F)\) vs \(\log|\varepsilon|\) gives the suppression order.
\item \textbf{R4 (MS-gate budget).} \(\text{MS}(\pi/2)=\exp(-i(\pi/4)(1+\varepsilon)XX)\), followed by two-qubit depolarizing channel of strength \(p_{2q}\). Average gate fidelity via the closed-form Haar-average \(F_{\text{avg}}=(dF_{\text{ent}}+1)/(d+1)\) with \(d=4\).
\item \textbf{RB-style (bonus).} \texttt{report/evidence/sim\_rb\_style.py} generates random single-qubit sequences from \(\{H,S,X_{\pi/2},X_{\pi}\}\), inverts each ideally, and fits an exponential to survival. Compares bare-\(H\) sequences to sequences that alternate \(H\) with \(H^{\dagger}\).
\end{enumerate}

\section{Results vs paper}
\begin{tabular}{p{0.6cm} p{6.5cm} p{3.5cm} p{2.5cm}}
\toprule
\# & Quantity & Paper & This run \\ \midrule
C3 & Fit 1-\(\sigma\) on \(\varepsilon,\phi,\delta\) & \(\sim 10^{-4}\) & \(5.4\times 10^{-4}\), \(8.0\times 10^{-4}\), \(1.3\times 10^{-5}\) \\
C4 & Suppression order, \(H\cdot H\) (bare) & \(O(\varepsilon^2)\) & \(|\varepsilon|^{1.99}\) \\
C4 & Suppression order, \(H\cdot H^{\dagger}\) (HI) & Cancels leading & \(1-F<10^{-15}\) (numerical zero) \\
C4 & Infidelity ratio bare/HI at \(\varepsilon=0.05\) & Large & \(7.5\times 10^{17}\) \\
C5 & \(F_{\text{MS}}\) at \(\varepsilon=0.09\), \(p_{2q}=0.02\) & \(97.5\%\) & \(F_{\text{avg}}=98.1\%\), \(F_{\text{ent}}=97.6\%\) \\
C6 & \(F_{\text{MS}}\) at eff.\ under-rot \(0.45\)~rad & \(\approx 91\%\) & \(F_{\text{avg}}=89.1\%\), \(F_{\text{ent}}=86.4\%\) \\
C7 & \(F_{\text{MS}}\) at \(\varepsilon=0.12\), \(p_{2q}=0.06\) & \(\approx 89\%\) & \(F_{\text{avg}}=94.8\%\), \(F_{\text{ent}}=93.5\%\) \\
--- & RB-style per-Clifford error, bare & --- & \(7.7\times 10^{-3}\) \\
--- & RB-style per-Clifford error, HI  & --- & \(5.2\times 10^{-3}\) \\
--- & Error-rate reduction (bare/HI)   & (qualitative) & \(1.49\times\) \\ \bottomrule
\end{tabular}

\medskip
\textbf{Interpretation:}
\begin{itemize}[leftmargin=1.4em]
\item C1--C4 are reproduced quantitatively. The core mathematical claim of the paper — that swapping $H$ for $H^\dagger$ cancels the coherent over-rotation error incurred by the adjacent gate — reproduces with a $\sim 10^{17}$-fold infidelity ratio at $\varepsilon=0.05$ (structurally exact under the paper's Eq.~1 noise model with $\phi=\delta=0$).
\item C3 (parameter recovery variance) reproduces to well within an order of magnitude of the paper's ``$\sim 10^{-4}$'' claim.
\item C5 and C6 reproduce within 0.5--2 percentage points on both average and entanglement fidelity conventions.
\item C7 is only partly reproduced (94.8\% vs 89\% target). The paper's reduced-cooling case includes a stochastic Debye--Waller reduction of the effective Rabi frequency due to $\bar n=0.5$ phonons; we used only the static over-rotation + depolarizing model, which is optimistic. Adding a stochastic Rabi-frequency distribution would degrade to the paper's number.
\item C8 (H$_2$ VQE end-to-end) was not attempted in this run — it is a system-level integration test, not a mechanism check.
\end{itemize}

\section{Verdict}
\textbf{REPLICATED (mechanism + budget).}
The Hidden-Inverses cancellation mechanism reproduces exactly under the paper's Eq.~(1) noise model, the SciPy noise-parameter fit reproduces the paper's variance order, and the MS-gate fidelity budget reproduces two of three regimes within experimental error and the third within the model-fidelity limits of our simplified stochastic treatment. The VQE H$_2$ system-level result was not attempted.

\section{Open Questions}
See \texttt{open\_questions.json} for the machine-readable version.
\begin{description}[leftmargin=1.4em]
\item[Q1] Under what noise-model conditions does HI \emph{fail}? Our $H\cdot H^{\dagger}$ test with $\phi=\delta=0$ gives numerical-zero infidelity. Rerunning with $\delta\ne 0$ or $\phi\ne 0$, or with time-dependent $\varepsilon(t)$ inside a single gate pulse, would tell us how brittle HI is compared to composite-pulse envelope-shaping methods that suppress errors to higher order in $\varepsilon$ regardless of drift model.
\item[Q2] For circuits where HI adjacent-cancellation cannot be arranged locally, can a global compiler pass find a set of decomposition choices that minimizes accumulated coherent error? The paper defers this ``open problem'' (Sec.~2) but does not give any complexity or heuristic. A branch-and-bound over $2^{|\text{self-adjoint gates}|}$ decompositions with the noise model as cost is finite for our tested lengths ($L\le 32$) and could be benchmarked.
\item[Q3] Our RB-style benchmark shows only 1.5$\times$ error reduction with alternating $H/H^{\dagger}$, well below the ratio suggested by the direct $H\cdot H$ vs $H\cdot H^{\dagger}$ test. The gap must come from non-adjacent Clifford insertions (S, X$_{\pi/2}$, X$_\pi$) breaking the HI cancellation. How does the effective HI benefit scale with the density of self-inverse gates in a given circuit family (e.g.\ QAOA on random graphs)?
\item[Q4] The paper reports MS-gate fidelity ``$97.5\%$'' as a single number derived from three independent physical parameters ($\bar n$, $\varepsilon$, $p_{2q}$). This is degenerate: many $(\bar n,\varepsilon,p_{2q})$ triples give the same fidelity. What identifiability test would separate cooling limits from laser-intensity miscalibration in future experiments, using the same amplification-circuit protocol used in Sec.~4?
\item[Q5] The paper compares HI to Randomized Compiling but not to explicit composite-pulse sequences (BB1, SK1, CORPSE). For small over-rotation $\varepsilon\ll 1$ our simulations show HI beats bare by many orders of magnitude when adjacent-cancellation works. In what parameter regime (gate count vs $\varepsilon$ magnitude vs drift timescale) do composite pulses beat HI, and where should a compiler prefer each?
\end{description}

\medskip
See also \texttt{workflow.md}, \texttt{artifacts\_summary.md}, and \texttt{failure\_analysis.md}.

\end{document}
